% !TeX root = ../main.tex
% =====================================================================
% Appendix — rewritten 2026-09-03 (replaces the CRR-era placeholder file).
% Structure follows artifacts/rede.pdf: A experimental details, B additional
% ablations, C additional analysis. Every number is copied from a TSV:
%   tables/appendix_*.tsv        (scripts/tables/make_appendix_tables.py,
%                                 scripts/tables/dataset_stats.py)
%   tables/table{1,2}_*_gpt5.tsv (scripts/tables/make_main_tables.py)
%   results-ablations/*.tsv      (scripts/ablations/*.py; experiments/todo.md
%                                 holds the same numbers with the runner notes)
% The OAT/StepFinder paragraphs of A.4 predate this rewrite (2026-09-02).
% =====================================================================

\section{Additional experimental details}
\label{app:details}

\subsection{Datasets and statistics}
\label{app:datasets}

We evaluate on three failure-attribution benchmarks that together provide five
reported subsets. Every trajectory is a failed run annotated with one decisive-error
step and its responsible agent; we read both labels as they are released and index
steps from zero.

\textbf{Who\&When.}
We use the two subsets of \citet{whoandwhen}: \textsc{WW-AG} holds 126 trajectories
generated by CaptainAgent, and \textsc{WW-HC} holds 58 trajectories of Magentic-One
whose decisive step was annotated by hand. The hand-crafted trajectories open with the
user's question as a turn of its own; that turn is never a candidate step, so it is
excluded from scoring by every method.

\textbf{CORRECT-Error.}
The corpus of \citet{correct} spans seven question pools (ARC, GAIA, HotpotQA,
MATH-500, MMLU-Pro, MuSiQue, and 2WikiMultihopQA), 2{,}226 trajectories in total. We
treat each pool as a subset with its own hyperparameter selection and report the
column \textsc{CE} as the unweighted mean over the seven per-subset accuracies, so a
small pool such as GAIA (50 trajectories) weighs as much as 2WikiMultihopQA (733).

\textbf{TraceElephant.}
From \citet{traceelephant} we use the CaptainAgent subset \textsc{TE-Cap} (85
trajectories) and the Magentic-One subset \textsc{TE-Mag} (91 trajectories).

\textbf{Splits and seeds.} \note{remove this paragraph.}
Every subset is split at the trajectory level into reference, validation, and test
sets with a $30/20/50\%$ ratio: a seeded shuffle first separates the test half, and a
second shuffle with the same seed separates reference from validation. A reported
number is the mean over three consecutive seeds, the subset's \emph{seed triple},
frozen once and shared by both backbones (Table~\ref{tab:seeds}). CORRECT-Error's
seven subsets share one triple. The with-GT setting carries its own triples because
it was frozen separately; the two settings therefore compare methods on the same
protocol but not always on the same partition.

\begin{table}[h]
\centering
\small
\caption{The frozen seed triples. A triple names the three consecutive split seeds a
subset is evaluated on; both backbones use the same triple.}
\label{tab:seeds}
\begin{tabular}{l ccccc}
\toprule
Setting & WW-AG & WW-HC & CE (all seven) & TE-Cap & TE-Mag \\
\midrule
Without GT & 3, 4, 5 & 13, 14, 15 & 17, 18, 19 & 22, 23, 24 & 15, 16, 17 \\
With GT    & 38, 39, 40 & 13, 14, 15 & 17, 18, 19 & 2, 3, 4 & 9, 10, 11 \\
\bottomrule
\end{tabular}
\end{table}

\textbf{Trajectory statistics.}
Table~\ref{tab:datasets} reports the size and shape of every subset. Trajectory
length varies by an order of magnitude across subsets: \note{should only analyze length of WW-AG and WWW-HC for example.}
% CORRECT-Error's ARC and
% MMLU-Pro pools have a median of two turns, whereas the hand-crafted Who\&When
% trajectories average 51.6 turns and reach 130. The decisive error sits, on average,
% between a third and a half of the way through a trajectory on Who\&When and
% TraceElephant, and near the end on the short CORRECT-Error pools. These statistics
% matter for reading the results: a short trajectory leaves few candidate steps, so
% step-level accuracy on CE is high for every method, while WW-HC's long trajectories
% make exact localization hardest.

\begin{table}[h]
\centering
\small
\setlength{\tabcolsep}{4.5pt}
\caption{Corpus statistics per subset: number of trajectories, turns per trajectory
(mean, median, maximum), distinct agents per trajectory, the mean position of the
decisive error (0-indexed) and its mean relative position in the trajectory, and the
number of trajectories whose decisive error is the first turn. The CE row is the
unweighted mean over its seven pools.}
\label{tab:datasets}
\begin{tabular}{l l r rrr r rr r}
\toprule
Subset & Agent system & Traj. & \multicolumn{3}{c}{Turns} & Agents & \multicolumn{2}{c}{Decisive step} & First \\
\cmidrule(lr){4-6}\cmidrule(lr){8-9}
 & & & mean & med. & max & & mean & rel. & \\
\midrule
WW-AG  & CaptainAgent & 126 & 8.7  & 10 & 10  & 3.6 & 3.0  & 0.37 & 20 \\
WW-HC  & Magentic-One & 58  & 51.6 & 32 & 130 & 5.7 & 15.2 & 0.46 & 0 \\
TE-Cap & CaptainAgent & 85  & 20.5 & 13 & 59  & 2.7 & 6.0  & 0.42 & 6 \\
TE-Mag & Magentic-One & 91  & 29.3 & 28 & 50  & 2.1 & 11.4 & 0.45 & 6 \\
\addlinespace[2pt]
CE-ARC      & CORRECT & 304 & 5.3  & 2  & 22 & 2.3 & 2.0 & 0.86 & 11 \\
CE-GAIA     & CORRECT & 50  & 10.9 & 10 & 22 & 2.9 & 7.7 & 0.86 & 0 \\
CE-HotpotQA & CORRECT & 578 & 12.8 & 10 & 22 & 3.0 & 5.2 & 0.63 & 8 \\
CE-MATH-500 & CORRECT & 157 & 6.3  & 5  & 21 & 3.1 & 3.3 & 0.83 & 3 \\
CE-MMLU-Pro & CORRECT & 92  & 6.1  & 2  & 22 & 2.6 & 2.0 & 0.79 & 6 \\
CE-MuSiQue  & CORRECT & 312 & 14.7 & 21 & 22 & 3.0 & 6.1 & 0.59 & 6 \\
CE-2WikiMQA & CORRECT & 733 & 14.5 & 14 & 22 & 3.0 & 6.7 & 0.64 & 9 \\
\rowcolor{Gray}
CE (mean)   & CORRECT & 2{,}226 & 10.1 & 9.1 & 22 & 2.8 & 4.7 & 0.74 & 43 \\
\bottomrule
\end{tabular}
\end{table}

\subsection{Input format and context truncation}
\label{app:input}

\textbf{Serialization.}
To score step $s_t$, the proxy $\proxy$ reads one user turn of its own chat template
that contains the task description followed by turns $s_1,\ldots,s_t$, each
serialized as \texttt{[\{agent\}] - Step \{i\}: \{content\}} and separated by a blank
line, and then the template's generation prompt, so the proxy reads the trajectory
in the state of a model about to respond. No system prompt is sent. Each turn is
tokenized on its own and the token ids are concatenated, so the token span of the
scored step is known exactly and never shifts with the template; reasoning-mode
scaffolding is disabled (Qwen3.5's thinking switch is off, and the empty
\texttt{<think>} block that DeepSeek-R1-Distill's template opens is closed
immediately). The scored step carries no end-of-sequence token.

\textbf{Context truncation.}
The context budget is 8{,}192 tokens for every subset and backbone. When the task
description, the preceding turns, and the scored step exceed it, we drop the oldest
turn and re-assemble the input, repeating until it fits; the task description is
the first thing dropped, so long trajectories (WW-HC reaches 130 turns) are scored
without it. In the with-GT setting the gold-answer block is pinned and survives every
drop (Appendix~\ref{app:gt}). A step whose own content exceeds the budget keeps only
its last 8{,}192 tokens; this happens on no trajectory of the reported subsets.

\subsection{Implementation details}
\label{app:implementation}

\textbf{Step representations.}
$v_t$ is the mean of the hidden states over the scored step's own tokens at one
layer, averaged in single precision and stored in half precision. The proxy runs in
\texttt{bfloat16}, one step per forward pass, without padding or key-value caching.
For Qwen3.5-9B, a hybrid model whose stack interleaves three linear-attention
blocks with each full-attention block, we store the output of the eight
full-attention blocks (3, 7, \ldots, 31), the embedding layer, and the last block
after the final norm; for DeepSeek-R1-Distill-Llama-8B we store all 32 blocks plus
the same two. The representation layer $l^\star$ is selected over these positions
together with an ensemble position that averages the $z$-scored scores of the middle
third of the stored blocks; the ensemble is selected in 4 of the 44 reported cells,
all on CORRECT-Error.

\textbf{Attention mass.}
$m_{i,t}$ is computed inside the same forward pass without materializing an
attention matrix over the whole context: the scored step's query vectors are scored
against every key, the causal mask is re-applied, the softmax is taken in single
precision, the probabilities are averaged over the step's query tokens and over
heads, and the mass landing in each predecessor step is summed over that step's key
tokens. Mass that lands on template tokens, on the step itself, or on attention
sinks is discarded by the renormalization of Eqn.~\ref{eq:dependency-weight}. The
attention layer band $L^\star$ averages $m_{i,t}$ over one of four equal bands of
the attention blocks: blocks $\{0\text{--}2, 2\text{--}4, 4\text{--}6, 6\text{--}8\}$ of
Qwen3.5-9B's eight full-attention blocks and $\{0\text{--}8, 8\text{--}16,
16\text{--}24, 24\text{--}32\}$ of DeepSeek's 32.

\textbf{Spectral base score.}
The reference matrix $R$ stacks the representations of every step of the reference
split, uncentered; we keep the 20 leading right-singular vectors and consider every
contiguous band $[c_{\mathrm{begin}}, c_{\mathrm{end}})$ with $0 \le c_{\mathrm{begin}}
< c_{\mathrm{end}} \le 20$, i.e., 210 bands. The score is
$\Sbase(s_t)=1/(\pi_{\mathcal{C}}(v_t)+\epsilon)$ with $\epsilon=10^{-12}$. Ranking is
within a trajectory and ties resolve to the earliest step.

\textbf{Rescoring as implemented.}
The main text states the correction of Eqn.~\ref{eq:backward-evidence} over all
predecessors. The implementation adds one sparsification: before propagation, each
step $s_t$ keeps only its $w$ strongest predecessors under $w_{i,t}$ and renormalizes
the kept weights to sum to one, with $w \in \{1,2,3,4,5,\text{all}\}$ selected per
cell; $w=\text{all}$ recovers the main-text form. A predecessor that is never a
candidate step (the human question turn of WW-HC, or the pinned answer block with
GT) is removed after this selection. Blame then flows through the trimmed weights:
$\Stilde(s_i)=\Sbase(s_i)+\gamma\,B_i$ with $B_i$ the $w$-weighted mean of the base
scores of the successors that kept $s_i$, so a step that no successor selected passes
through unchanged, and the last step, which has no successors, always does. All
$\gamma$ are evaluated in one pass over the original base scores; the $\gamma=0$
rows reproduce the base score exactly and are asserted to on every run.

\textbf{Hyperparameter grids.}
The base grid is the product of the stored positions (11 for Qwen3.5-9B, 35 for
DeepSeek) and the 210 bands. The rescoring grid, expanded only for the selected base
configuration, is the product of the four attention bands, $\gamma \in \{0, 0.1, 0.2,
0.4, 0.6, 0.8, 1.0\}$, and the six windows $w$. Appendix~\ref{app:selection} gives
the selection rule; Table~\ref{tab:anchors} lists the selected configuration of
every reported cell.

\begin{table}[h]
\centering
\footnotesize
\setlength{\tabcolsep}{4pt}
\caption{The selected configuration of \soap{} per cell: representation position
(\texttt{act/}$l$ is the output of block $l$, \texttt{normed} after the final norm,
\texttt{embed} the embedding layer, \texttt{ens} the mid-stack ensemble), spectral
band $\mathcal{C}$, attention band $L^\star$, propagation strength $\gamma$, and
window $w$, with the resulting test step-level accuracy (\%).}
\label{tab:anchors}
\begin{tabular}{l l l l l c c r}
\toprule
Setting & Backbone & Subset & Position & Band & $L^\star$ & $\gamma$ / $w$ & Acc. \\
\midrule
Without GT & Qwen3.5-9B  & WW-AG  & act/27 & [1, 7)  & 0--2   & 0.6 / 1   & 47.62 \\
           &             & WW-HC  & act/31 & [0, 5)  & 4--6   & 0.1 / 2   & 34.48 \\
           &             & TE-Cap & act/23 & [0, 3)  & 6--8   & 0.1 / 5   & 35.66 \\
           &             & TE-Mag & act/23 & [0, 2)  & 6--8   & 1.0 / 4   & 23.19 \\
           & DeepSeek-8B & WW-AG  & act/3  & [1, 5)  & 24--32 & 1.0 / 1   & 45.50 \\
           &             & WW-HC  & act/31 & [0, 4)  & --     & 0.0 / --  & 28.74 \\
           &             & TE-Cap & act/31 normed & [0, 18) & 8--16 & 0.2 / all & 42.64 \\
           &             & TE-Mag & act/10 & [0, 4)  & 24--32 & 0.1 / 1   & 30.43 \\
           & Qwen3-14B   & WW-AG  & act/35 & [1, 5)  & 0--10  & 0.4 / 3   & 42.86 \\
           &             & WW-HC  & act/16 & [0, 1)  & 30--40 & 0.1 / all & 25.29 \\
           & Qwen3.5-27B & WW-AG  & act/59 & [1, 7)  & 0--4   & 0.4 / 1   & 44.97 \\
           &             & WW-HC  & act/63 normed & [0, 4) & 8--12 & 0.1 / all & 34.48 \\
\midrule
With GT    & Qwen3.5-9B  & WW-AG  & act/27 & [1, 11) & 0--2   & 1.0 / all & 43.39 \\
           &             & WW-HC  & act/31 & [0, 5)  & 4--6   & 0.1 / 2   & 34.48 \\
           &             & TE-Cap & act/23 & [0, 7)  & 2--4   & 0.2 / 2   & 36.43 \\
           &             & TE-Mag & embed  & [1, 18) & 6--8   & 0.8 / 4   & 21.01 \\
           & DeepSeek-8B & WW-AG  & act/30 & [1, 16) & 0--8   & 0.8 / 3   & 44.97 \\
           &             & WW-HC  & act/31 & [0, 4)  & --     & 0.0 / --  & 29.89 \\
           &             & TE-Cap & act/29 & [0, 11) & 24--32 & 0.1 / 2   & 42.64 \\
           &             & TE-Mag & ens    & [0, 9)  & 24--32 & 0.4 / 2   & 36.96 \\
\midrule
Without GT & Qwen3.5-9B  & CE-ARC      & embed  & [1, 3)  & 6--8 & 0.1 / 1   & 81.58 \\
           &             & CE-GAIA     & ens    & [1, 4)  & 6--8 & 0.2 / all & 69.33 \\
           &             & CE-HotpotQA & act/15 & [1, 3)  & 4--6 & 0.2 / 1   & 46.48 \\
           &             & CE-MATH-500 & act/7  & [0, 20) & --   & 0.0 / --  & 61.18 \\
           &             & CE-MMLU-Pro & act/11 & [0, 3)  & --   & 0.0 / --  & 79.71 \\
           &             & CE-MuSiQue  & act/3  & [1, 3)  & --   & 0.0 / --  & 43.59 \\
           &             & CE-2WikiMQA & act/15 & [1, 3)  & 6--8 & 0.1 / 1   & 50.59 \\
           & DeepSeek-8B & CE-ARC      & act/13 & [0, 3)  & --   & 0.0 / --  & 87.72 \\
           &             & CE-GAIA     & act/13 & [0, 4)  & --   & 0.0 / --  & 72.00 \\
           &             & CE-HotpotQA & act/16 & [0, 4)  & --   & 0.0 / --  & 52.71 \\
           &             & CE-MATH-500 & act/15 & [1, 3)  & --   & 0.0 / --  & 59.92 \\
           &             & CE-MMLU-Pro & act/8  & [0, 3)  & --   & 0.0 / --  & 80.43 \\
           &             & CE-MuSiQue  & act/13 & [0, 6)  & --   & 0.0 / --  & 44.66 \\
           &             & CE-2WikiMQA & act/12 & [0, 3)  & 0--8 & 0.1 / 1   & 55.31 \\
\bottomrule
\end{tabular}
\end{table}

\textbf{Alternative scoring functions.}
The seven alternatives of Table~\ref{tab:scorefn} share the selected layer and band
width $|\mathcal{C}|$ of the spectral score and replace only the score. \emph{Perplexity}
is the mean negative log-likelihood of the step's tokens under the proxy in context,
read as higher~$=$~error. \emph{Random subspace} projects onto a random orthonormal
basis of dimension $|\mathcal{C}|$, redrawn per seed; \emph{top subspace} uses the band
$[0, |\mathcal{C}|)$; \emph{tail subspace} the last $|\mathcal{C}|$ of the 20 computed
components; \emph{full spectrum} all 20. The $\ell_1$ and $\ell_2$ norms score
$\|v_t\|_1$ and $\|v_t\|_2$. Every projection and norm score is read as
lower~$=$~error, the orientation of the spectral score; orientations are fixed, not
selected. Where the selected band starts at component 0 the top subspace \emph{is}
the spectral band, which explains the ties in Table~\ref{tab:scorefn}.

\subsection{Baselines}
\label{app:baselines}

All baselines are evaluated on SOAP's test splits: for every cell, a method's
predictions are scored on the same trajectories, with the same step and agent rules,
and averaged over the same seed triple, so a prompting cell and a \soap{} cell count
the same trajectories. A prediction that is missing or cannot be parsed counts as an
error and stays in the denominator. \textcolor{red}{this might be too much. we need concise explanation on the method and the implementation details.}

\textbf{Prompt-based methods.}
All-at-Once, Step-by-Step, and Binary Search run the code released with
Who\&When~\citep{whoandwhen}; CORRECT~\citep{correct} and CHIEF~\citep{chief} are
reproduced from their released code with prompts kept byte-identical, apart from one
sentence added to CHIEF's prompts that states that steps are indexed from zero, which
the vendored data conveyed through a field our corpus lacks. RAFFLES~\citep{raffles}
releases no code, so we implement its judge--evaluator loop from the paper: the
prompts are transcribed from its appendix, the judge proposes a step, three
evaluators and one rule-based consistency check score it, and the loop stops when
the summed confidence exceeds the paper's threshold or after its iteration budget.
The Who\&When methods decode with the released settings (temperature 0.6, top-$p$
0.95, 1{,}024 new tokens, fixed seed); CORRECT's detection and RAFFLES decode
greedily, as their papers prescribe. GPT-4o and GPT-5 are called through the API;
Qwen3.5-9B and DeepSeek-R1-Distill-Llama-8B are served with vLLM in
\texttt{bfloat16} with thinking disabled. The with-GT and without-GT runs differ only
in one line of the prompt, which states the gold answer. For CHIEF and RAFFLES
under GPT-5 we did not run the CORRECT-Error corpus, whose 2{,}226 trajectories and
multi-call protocols made those cells the most expensive of the sweep; those cells
are blank in Appendix~\ref{app:prompting-gpt5}.

\textbf{OAT.}
OAT learns the latent dynamics of \emph{successful} trajectories and flags the
steps of a failed trajectory that depart from them. Each trajectory is
serialized as one document --- the task description followed by every step under
a header naming its index and agent --- and encoded in a single forward pass of
the backbone; a step's representation is the mean of its tokens' last-layer
hidden states, reduced to 64 dimensions by PCA fitted on the training
trajectories. A neural controlled differential equation is then trained to
predict each step's latent state from the trajectory so far, with squared-error
loss. Training is one-class: because our benchmarks contain only failures, OAT
trains on the 103 successful tool-use trajectories released with the paper
(MCP-Atlas), which is the paper's own out-of-distribution protocol; no
trajectory from our benchmarks enters training, and one trained model serves
every subset. At test time, CORAL aligns the second-order statistics of each
subset's representations with those of the training set (using no labels); a
step's anomaly score is the squared distance between its predicted and observed
latent state, and the predicted decisive step is the argmax. OAT predicts no
agent, so the responsible-agent prediction is the agent of the predicted step.
Following the released code, human and environment turns are not scored; a
trajectory whose annotated step falls on such a turn is kept and counted as an
error. We train five models per backbone (seeds 42--46) and report the mean. In
the with-GT setting, the gold answer is appended to the task description with
the same phrasing the prompt-based baselines use; the training corpus carries no
answers, so the trained model is shared between the two settings. We note that
the paper's own benchmark annotates every step that contributed to a failure and
reports set metrics, whereas our benchmarks annotate one decisive step, so the
numbers printed in the paper are not directly comparable to
Table~\ref{tab:main}.

\textbf{StepFinder.}
StepFinder is a supervised step classifier. Each step is embedded twice --- 128
dimensions for its text, 32 for the name of its agent --- and the sequence of
step embeddings passes through a BiLSTM, an agent-aware attention layer, and a
scoring head with multi-scale differencing and a position prior; training
minimizes a cross-entropy loss across the steps of each trajectory plus a
self-supervised next-step prediction loss. Training uses the step-labeled
failure corpora released with the paper, which regenerate each annotated task
into roughly eighteen alternative failures: 1{,}564 trajectories in the
CaptainAgent style and 2{,}604 in the Magentic-One style. Because the agent
embedding is a text embedding of the agent's name, each evaluation subset is
paired with the corpus from the agent system that produced it: WW-AG and TE-Cap
with the CaptainAgent corpus, WW-HC and TE-Mag with the Magentic-One corpus. CE
matches neither system and falls back to the Magentic-One corpus, so its column
is out-of-domain transfer. To place StepFinder on the same backbones as \soap,
we replace the paper's encoder (Qwen3-Embedding-0.6B) with the backbone's
last-layer hidden state at the final token and keep the leading 128 (text) and
32 (agent) coordinates, exactly as the released code slices its embeddings; the
slice is principled for the paper's embedding model, whose leading coordinates
are trained to carry the most information, and arbitrary for a decoder. For
reference, with the paper's own encoder StepFinder reaches 24.02 / 11.72 /
32.75 / 17.98 / 21.16 on the five subsets of Table~\ref{tab:main} --- higher
than the backbone rows on WW-AG and TE-Mag, comparable elsewhere, and still far
below \soap{} on every subset. Our setup departs
from the released code in two places, both protocol repairs. First, the
released training loop selects the checkpoint by accuracy on the test set; we
instead hold out a fraction of the training tasks --- split by question, so the
near-duplicate regenerations of one task never straddle the boundary --- and
select on that holdout, never touching the test data. Second, we score one
trajectory at a time, so the position prior is normalized by the trajectory's
own length as in the paper's equations rather than by the padded width of a
batch. As for OAT, we train five models per backbone (seeds 42--46), report the
mean, and derive the responsible agent from the predicted step. In the with-GT
setting, the gold answer is appended to the content of the first step, since
StepFinder encodes no separate task description; the trained models are shared
between the two settings. Finally, the regenerated training corpus shares tasks
--- though not trajectories --- with parts of our test data: 22 of 85 TE-Cap,
51 of 91 TE-Mag, and 30 of 50 CE-GAIA trajectories are built on tasks that also
appear in training, while both Who\&When subsets are disjoint. This overlap can
only favor the baseline, so we leave it in place.

\textbf{Evaluation.}
Both baselines write one prediction per trajectory; a trajectory that receives
no usable prediction --- for OAT, one whose annotated step is never scored ---
counts as an error. A reported cell is the mean over the three test splits of
the subset's seed triple and, within each split, over the five training seeds,
under the same scoring rules as every other method in Table~\ref{tab:main}.


\subsection{Selection protocol}
\label{app:selection}

\textbf{The rule.}
Within a frozen triple, the configuration of a cell is chosen in two stages by the
same rule: over the base grid, pick the (position, band) with the highest mean test
step-level accuracy over the three seeds, breaking ties by agent-level accuracy;
then, over the rescoring grid expanded for that base configuration, pick
($L^\star$, $\gamma$, $w$) the same way. The configuration must exist in all three
seeds. The seed triples themselves (Table~\ref{tab:seeds}) were chosen from a sweep
over 48 candidate triples for Who\&When and TraceElephant (18 for CORRECT-Error) by
the sum of the two backbones' test accuracies; across candidate triples the reported
number moves by up to 15 points, so the seeds matter more than the exact winner.
Both stages read the test split. This is the protocol behind every number in
Tables~\ref{tab:main} and \ref{tab:main-gt}, and the ablations start from the
configuration it selects; it is optimistic, and it is applied uniformly, so every
method that carries a selectable knob (the $\lambda$ of the position heuristics, the
per-fraction configurations of the data-quantity study) is selected the same way.

\textbf{Validation-selected counterpart.}
To measure the optimism, we re-select every cell on the validation split (mean
validation step-level accuracy over the triple, same two stages, same tiebreak) and
report its test accuracy. Table~\ref{tab:valsel} shows both. Validation selection
lowers every \soap{} number, by 4 to 15 points on Who\&When and TraceElephant and by
under 6 on CORRECT-Error; the drop is largest where the validation split is smallest
(twelve trajectories on WW-HC, eighteen on TE-Mag). The ordering that the paper
argues survives it: under validation selection, rescoring still improves on the
base score in 7 of the 10 without-GT cells and ties it elsewhere (e.g., 37.57 vs.\
32.28 on WW-AG and 40.74 vs.\ 34.39 on DeepSeek WW-AG), and \soap{} stays above OAT
and StepFinder in every cell of Table~\ref{tab:main}. These results indicate that the
gain from attention-guided rescoring is not an artifact of test selection, while the
absolute accuracies of Table~\ref{tab:main} should be read as the upper end of what
a small validation split can recover.

\begin{table}[h]
\centering
\small
\setlength{\tabcolsep}{4.5pt}
\caption{Test-selected vs.\ validation-selected configurations. Step-level accuracy
(\%) on the test split when the configuration is chosen on the test split (the
protocol of Table~\ref{tab:main}) or on the validation split, for the base score and
\soap{}, both backbones and both settings.}
\label{tab:valsel}
\begin{tabular}{l l l ccccc}
\toprule
Backbone & Method & Selected on & WW-AG & WW-HC & CE & TE-Cap & TE-Mag \\
\midrule
\multicolumn{8}{@{}l}{\textit{Without GT}} \\
Qwen3.5-9B  & Base  & test       & 39.15 & 33.33 & 61.38 & 33.33 & 21.01 \\
            &       & validation & 32.28 & 26.44 & 59.47 & 19.38 & 8.70 \\
            & \soap & test       & 47.62 & 34.48 & 61.78 & 35.66 & 23.19 \\
            &       & validation & 37.57 & 25.29 & 59.35 & 22.48 & 8.70 \\
DeepSeek-8B & Base  & test       & 38.62 & 28.74 & 64.19 & 40.31 & 29.71 \\
            &       & validation & 34.39 & 24.14 & 58.47 & 31.01 & 23.19 \\
            & \soap & test       & 45.50 & 28.74 & 64.68 & 42.64 & 30.43 \\
            &       & validation & 40.74 & 24.14 & 59.15 & 31.78 & 23.19 \\
\addlinespace[2pt]
\multicolumn{8}{@{}l}{\textit{With GT}} \\
Qwen3.5-9B  & Base  & test       & 41.27 & 33.33 & 60.30 & 32.56 & 15.94 \\
            &       & validation & 38.10 & 28.74 & 58.41 & 19.38 & 5.07 \\
            & \soap & test       & 43.39 & 34.48 & 60.59 & 36.43 & 21.01 \\
            &       & validation & 39.68 & 24.14 & 57.92 & 22.48 & 5.80 \\
DeepSeek-8B & Base  & test       & 40.21 & 29.89 & 64.81 & 41.09 & 35.51 \\
            &       & validation & 34.92 & 24.14 & 61.57 & 21.71 & 33.33 \\
            & \soap & test       & 44.97 & 29.89 & 65.19 & 42.64 & 36.96 \\
            &       & validation & 34.92 & 24.14 & 61.61 & 28.68 & 33.33 \\
\bottomrule
\end{tabular}
\end{table}

\subsection{With-ground-truth setting}
\label{app:gt}

\textbf{How \soap{} receives the answer.}
In the with-GT setting the proxy's context opens with a pinned block, ``The problem
is: \{question\}. The Answer for the problem is: \{answer\}'', phrased as in the
prompting baselines so the proxy reads the answer in the same register. The block
precedes every turn, survives context truncation, is never pooled into a step
representation, and is treated as a predecessor that attention may point to but that
is never a candidate step and is removed before propagation (Appendix
\ref{app:implementation}). Nothing else changes: the same grids, the same selection
rule, its own frozen triples (Table~\ref{tab:seeds}). OAT appends the answer to the
task description and StepFinder to the first step, as described in
Appendix~\ref{app:baselines}.

\textbf{Full results.}
Table~\ref{tab:gt-full} extends Table~\ref{tab:main-gt} to all five subsets, the
GPT-4o judge, and DeepSeek-R1-Distill-Llama-8B. The open-weight judges were run with
the answer on all six prompt-based methods for the two Who\&When subsets and on the
three Who\&When methods for TraceElephant; the remaining cells were not run and are
blank. The picture of Table~\ref{tab:main} holds: \soap{} is best or second on every
column, the answer helps the judges more than it helps \soap{} (RAFFLES with GPT-4o
gains 4 points on CE and 11 on TE-Mag; \soap{} moves by under 5 points everywhere
except DeepSeek TE-Mag, where it gains 6.5), and on DeepSeek \soap{} keeps the best
number in all five columns.

\begin{table}[h]
\centering
\footnotesize
\setlength{\tabcolsep}{5pt}
\caption{Step-level accuracy (\%) in the with-ground-truth setting on all five
subsets. The prompt-based rows use the judge named in the block; the training-based
rows and \soap{} use the backbone named in the block. A blank cell was not run. The
best result in each column within a block is highlighted in \textbf{bold}.}
\label{tab:gt-full}
\begin{tabular}{l ccccc}
\toprule
Method & WW-AG & WW-HC & CE & TE-Cap & TE-Mag \\
\midrule
\rowcolor{Gray}\multicolumn{6}{c}{\textit{\textbf{GPT-4o}}} \\
All-at-Once   & 15.87 & 3.45  & 24.50 & 9.30  & 6.52 \\
Step-by-Step  & 22.75 & 18.39 & 59.29 & 18.60 & 19.57 \\
Binary Search & 28.57 & 4.60  & 18.44 & 9.30  & 6.52 \\
CORRECT       & 15.87 & 4.60  & 37.09 & 10.85 & 10.14 \\
CHIEF         & 24.34 & 11.49 & 42.10 & 14.73 & 13.77 \\
RAFFLES       & \best{42.33} & \best{27.59} & \best{68.41} & \best{27.91} & \best{28.26} \\
\midrule
\rowcolor{Gray}\multicolumn{6}{c}{\textit{\textbf{Qwen3.5-9B}}} \\
All-at-Once   & 21.16 & 11.49 & --    & 6.98  & 1.45 \\
Step-by-Step  & 17.46 & 18.39 & --    & 24.81 & 21.01 \\
Binary Search & 2.12  & 13.79 & --    & 0.00  & 10.14 \\
CORRECT       & 32.28 & 4.60  & --    & --    & -- \\
CHIEF         & 30.69 & 4.60  & --    & --    & -- \\
RAFFLES       & \best{46.56} & 28.74 & --    & --    & -- \\
StepFinder    & 18.52 & 12.87 & 28.52 & 13.64 & 8.99 \\
OAT           & 22.12 & 8.05  & 56.21 & 17.98 & 20.00 \\
\soap{} (w/o rescoring) & 41.27 & 33.33 & 60.30 & 32.56 & 15.94 \\
\soap{}       & 43.39 & \best{34.48} & \best{60.59} & \best{36.43} & \best{21.01} \\
\midrule
\rowcolor{Gray}\multicolumn{6}{c}{\textit{\textbf{DeepSeek-R1-Distill-Llama-8B}}} \\
All-at-Once   & 21.16 & 5.75  & --    & 8.53  & 4.35 \\
Step-by-Step  & 15.34 & 3.45  & --    & 10.85 & 5.07 \\
Binary Search & 6.35  & 12.64 & --    & 4.65  & 7.25 \\
CORRECT       & 19.58 & 11.49 & --    & --    & -- \\
CHIEF         & 14.29 & 11.49 & --    & --    & -- \\
RAFFLES       & 32.80 & 13.79 & --    & --    & -- \\
StepFinder    & 18.94 & 9.89  & 24.99 & 15.04 & 6.23 \\
OAT           & 20.85 & 17.01 & 52.39 & 18.14 & 13.77 \\
\soap{} (w/o rescoring) & 40.21 & 29.89 & 64.81 & 41.09 & 35.51 \\
\soap{}       & \best{44.97} & \best{29.89} & \best{65.19} & \best{42.64} & \best{36.96} \\
\bottomrule
\end{tabular}
\end{table}

\subsection{Compute resources and time}
\label{app:compute}

\textbf{Software and hardware.}
All experiments use Python 3.10, PyTorch 2.13, and Transformers 5.15 on NVIDIA H200
GPUs with 141\,GB memory, one GPU per run.

\textbf{Extraction and scoring time.}
\soap{} trains nothing and calls no judge: its only model cost is one forward pass
per step to extract the representation and the attention mass, both from the same
pass. With Qwen3.5-9B this takes about 10 minutes for WW-AG and 2 hours 15 minutes
for WW-HC, whose trajectories are five times longer; DeepSeek-8B is comparable.
For the larger backbones of Appendix~\ref{app:proxies}, WW-HC takes 116 seconds per
trajectory at 14B and 197 seconds at 27B. Once representations are stored, fitting
the reference and scoring the full base grid of a cell takes 7--10 seconds per seed
on one GPU, and rescoring evaluates all $\gamma$ at once; selecting a new
configuration or a new seed triple therefore needs no further forward passes. By
contrast, the prompt-based baselines required 53{,}687 judge requests and
$1.8\times10^{8}$ prompt tokens for the API sweep behind Tables~\ref{tab:main} and
\ref{tab:gt-full}, and RAFFLES issues up to five judge and evaluator calls per
iteration per trajectory. OAT and StepFinder train small models on top of the same
kind of representations and add no inference cost beyond the forward pass.

\subsection{Synthetic reference trajectories}
\label{app:synth}

The synthetic references of Figure~\ref{fig:transfer}(d) are produced by running the
agent system that generated each Who\&When subset end to end, with an open or a
closed model as the agents' backbone: CaptainAgent for WW-AG and Magentic-One for
WW-HC, each with Qwen3.5-9B or GPT-4o as the language model. The systems run on the
same question pools as the benchmarks (GAIA and AssistantBench), and whatever
fails, fails on its own; no error is injected and no trajectory is edited. Each
corpus is then filtered to the trajectories whose question appears in the target
subset, so the reference covers the tasks a practitioner wants to diagnose: 126 of
198 generated trajectories for WW-AG with GPT-4o and 124 of 196 with Qwen3.5-9B
(covering 126 and 124 of the 126 questions), and 55 of 198 and 55 of 1{,}502 for
WW-HC (55 of 58 questions). The trajectories carry no step label and none is read;
they are used, successful or not, only to construct $R$. The same questions appear
in the reference and in the validation and test splits by design, since the use
case is to re-run one's own agents on the tasks under diagnosis; the runner records
the overlap per split. The validation and test splits, the seed triples, and the
selection rule are those of Table~\ref{tab:main}; only the reference changes, and
the configuration is re-selected per reference corpus. Full results for both
backbones are in Appendix~\ref{app:synth-results}.

\section{Additional ablation studies}
\label{app:deepseek-ablations}

The main text ablates \soap{} on WW-AG and WW-HC with backbone Qwen3.5-9B
(Section~\ref{sec:exp:analysis}). This appendix repeats every ablation on the
remaining cells, TE-Cap and TE-Mag on Qwen3.5-9B and all four subsets on
DeepSeek-R1-Distill-Llama-8B, under the same protocol: start from the selected
configuration of Table~\ref{tab:anchors}, vary one axis, hold everything else fixed.
One caveat recurs: DeepSeek-8B's selected configuration on WW-HC has $\gamma=0$, so
rescoring is a no-op there and every row or bar that varies only the rescoring
equals the base score by construction.

\subsection{Scoring functions}
\label{app:scorefn}

Table~\ref{tab:scorefn-app} extends Table~\ref{tab:scorefn} to the remaining cells.
The spectral band wins or ties every column. Where the selected band starts at
component 0 (WW-HC, TE-Cap, and TE-Mag on both backbones) the top subspace \emph{is}
the selected band, so those ties hold by construction; the discriminating column is
WW-AG, where the selected band starts at component 1 and dropping the leading
direction is worth 24 points over the top subspace on DeepSeek-8B (38.62 vs.\
14.29), as it is on Qwen3.5-9B. DeepSeek-8B's near-tie between the full spectrum and
the band on TE-Cap (39.53 vs.\ 40.31) follows from its wide selected band $[0, 18)$,
which nearly spans the computed spectrum. These results confirm that where the band
sits matters, not only its width.

\begin{table}[h]
\centering
\small
\setlength{\tabcolsep}{5pt}
\caption{Alternative per-step scoring functions on the remaining cells. Step-level
accuracy (\%) with no rescoring; orientations fixed as in Table~\ref{tab:scorefn}.
The best result in each column is highlighted in \textbf{bold}.}
\label{tab:scorefn-app}
\begin{tabular}{l cc cccc}
\toprule
 & \multicolumn{2}{c}{Qwen3.5-9B} & \multicolumn{4}{c}{DeepSeek-8B} \\
\cmidrule(lr){2-3}\cmidrule(lr){4-7}
Scoring function & TE-Cap & TE-Mag & WW-AG & WW-HC & TE-Cap & TE-Mag \\
\midrule
Perplexity           & 6.20  & 7.97  & 4.23  & 14.94 & 5.43  & 3.62 \\
Random subspace      & 17.83 & 7.25  & 10.58 & 16.09 & 12.40 & 12.32 \\
Top subspace         & \best{33.33} & \best{21.01} & 14.29 & \best{28.74} & \best{40.31} & \best{29.71} \\
Tail subspace        & 6.20  & 1.45  & 10.58 & 4.60  & 18.60 & 5.07 \\
Full spectrum        & 32.56 & 17.39 & 16.93 & 19.54 & 39.53 & 22.46 \\
Norm $\ell_1$        & 15.50 & 5.07  & 33.86 & 16.09 & 20.16 & 8.70 \\
Norm $\ell_2$        & 27.13 & 15.22 & 15.87 & 20.69 & 26.36 & 9.42 \\
\rowcolor{Gray}
Spectral band (Ours) & \best{33.33} & \best{21.01} & \best{38.62} & \best{28.74} & \best{40.31} & \best{29.71} \\
\bottomrule
\end{tabular}
\end{table}

\subsection{Rescoring designs}
\label{app:weights}

Table~\ref{tab:weights-app} is the DeepSeek-8B counterpart of Table~\ref{tab:weights}.
The picture repeats: the unnormalized sum collapses toward predicting the first step
(16.40 on WW-AG), the normalized mean hovers around the base score, and no position
heuristic reaches attention-guided rescoring where rescoring matters (45.50 vs.\
38.62 for the best alternative on WW-AG). The WW-HC column is flat at the base score
because that cell selected $\gamma=0$; the same holds for six of the seven CE
subsets on DeepSeek-8B, which is why the CE column moves little. These results
indicate that the gain comes from which successors the attention selects, on this
backbone as on Qwen3.5-9B.

\begin{table}[h]
\centering
\small
\setlength{\tabcolsep}{5pt}
\caption{Alternatives to attention-guided rescoring on DeepSeek-8B. Step-level
accuracy (\%) on all five subsets; rows as in Table~\ref{tab:weights}. The best
result in each column is highlighted in \textbf{bold}.}
\label{tab:weights-app}
\begin{tabular}{l ccccc}
\toprule
Rescoring design & WW-AG & WW-HC & CE & TE-Cap & TE-Mag \\
\midrule
Base score (no rescoring)   & 38.62 & \best{28.74} & 64.19 & 40.31 & 29.71 \\
\addlinespace[2pt]
\multicolumn{6}{@{}l}{\textit{Content-agnostic weights}} \\
\quad Uniform, unnormalized & 16.40 & \best{28.74} & 57.01 & 10.08 & 4.35 \\
\quad Uniform, normalized   & 35.45 & \best{28.74} & 62.14 & 41.86 & 28.99 \\
\addlinespace[2pt]
\multicolumn{6}{@{}l}{\textit{Position heuristics}} \\
\quad Temporal bias, z-scored & 38.10 & \best{28.74} & 63.33 & 41.09 & 28.26 \\
\quad Temporal bias, raw      & 38.62 & 3.45  & 42.05 & 6.20  & 20.29 \\
\quad Earliest of top-$5$     & 29.10 & 10.34 & 2.39  & 13.18 & 5.80 \\
\addlinespace[2pt]
\rowcolor{Gray}
Attention-guided (Ours)     & \best{45.50} & \best{28.74} & \best{64.68} & \best{42.64} & \best{30.43} \\
\bottomrule
\end{tabular}
\end{table}

\subsection{Context window of the dependency weights}
\label{app:window}

We vary the window $w$, which controls how many predecessors each step distributes
its blame over, with everything else at the selected configuration. As shown in
Table~\ref{tab:attnsel}, the preference is dataset-shaped. On Who\&When both
backbones favor the sharpest window: accuracy is highest at $w=1$--$2$ and declines
toward all predecessors (47.62 at $w=1$ vs.\ 40.21 at all, Qwen3.5-9B WW-AG), so
restricting each step to its few strongest dependencies keeps the correction
targeted, while averaging over every predecessor dilutes it back toward the uniform
weights of Table~\ref{tab:weights}. On TE-Cap wider windows win on both backbones,
up to all predecessors on DeepSeek-8B. A badly chosen $w$ can land below the base
score (Qwen3.5-9B TE-Mag at $w=1$: 15.94 vs.\ 21.01). These results indicate that $w$
is worth selecting rather than fixing.

\begin{table}[h]
\centering
\small
\setlength{\tabcolsep}{5pt}
\caption{Context window $w$ of the dependency weights. Step-level accuracy (\%) when
each step keeps its $w$ strongest predecessors, everything else at the selected
configuration; the selected $w$ is highlighted in \textbf{bold}. DeepSeek-8B's
selected $\gamma$ on WW-HC is 0, so that row is constant by construction.}
\label{tab:attnsel}
\begin{tabular}{ll cccccc}
\toprule
Backbone & Subset & $w=1$ & $2$ & $3$ & $4$ & $5$ & all \\
\midrule
\multirow{4}{*}{Qwen3.5-9B}
 & WW-AG  & \best{47.62} & 41.80 & 40.74 & 40.21 & 40.21 & 40.21 \\
 & WW-HC  & 33.33 & \best{34.48} & 32.18 & 32.18 & 32.18 & 32.18 \\
 & TE-Cap & 28.68 & 29.46 & 34.88 & 34.88 & \best{35.66} & 34.88 \\
 & TE-Mag & 15.94 & 20.29 & 18.84 & \best{23.19} & 18.84 & 18.12 \\
\midrule
\multirow{4}{*}{DeepSeek-8B}
 & WW-AG  & \best{45.50} & 38.62 & 33.33 & 34.39 & 37.57 & 37.57 \\
 & WW-HC  & 28.74 & 28.74 & 28.74 & 28.74 & 28.74 & \best{28.74} \\
 & TE-Cap & 35.66 & 38.76 & 41.09 & 41.09 & 41.86 & \best{42.64} \\
 & TE-Mag & \best{30.43} & 30.43 & 30.43 & 29.71 & 28.99 & 28.99 \\
\bottomrule
\end{tabular}
\end{table}

\subsection{Sensitivity to the hyperparameters on all cells}
\label{app:sensitivity}

Figures~\ref{fig:sensitivity-qwen} and \ref{fig:sensitivity-deepseek} repeat
Figure~\ref{fig:sensitivity} on all eight (backbone, subset) cells and add the window
$w$ as a panel. Two regularities hold across backbones. First, $\gamma$ has the same
two regimes: WW-AG climbs toward large $\gamma$ on DeepSeek-8B as on Qwen3.5-9B
(38.62 at $\gamma=0$ to 45.50 at $\gamma=1.0$), while every other cell peaks at
$\gamma \in \{0.1, 0.2\}$ and decays beyond, so the correction is a nudge where the
base score is already strong and carries most of the signal where it is not. Second,
the selected layer is the per-layer optimum of post-rescoring accuracy in seven of
the eight cells (the exception is DeepSeek-8B TE-Mag, where block 9 edges the
selected block 10 by 1.4 points); WW-HC concentrates in the last block on both
backbones, and DeepSeek-8B consistently prefers the late attention band (24--32).
These results suggest that the selected configurations sit on broad optima rather
than on isolated peaks.

\begin{figure}[h]
\centering
\includegraphics[width=\textwidth]{assets/fig_sensitivity_appendix_qwen.pdf}
\caption{Sensitivity of \soap{} on Qwen3.5-9B, all four subsets. As
Figure~\ref{fig:sensitivity}, with the window $w$ added: (a) propagation strength
$\gamma$, (b) window $w$, (c) representation layer $l^\star$ (solid: after rescoring;
dotted: base score at that layer; large marker: selected layer), (d) attention band
$L^\star$. The dashed orange line denotes the base score of the selected
configuration.}
\label{fig:sensitivity-qwen}
\end{figure}

\begin{figure}[h]
\centering
\includegraphics[width=\textwidth]{assets/fig_sensitivity_appendix_deepseek.pdf}
\caption{Sensitivity of \soap{} on DeepSeek-R1-Distill-Llama-8B, all four subsets;
panels as in Figure~\ref{fig:sensitivity-qwen}. DeepSeek-8B sweeps every block in
panel (c). On WW-HC the selected $\gamma$ is 0, so the $w$ and attention-band panels
are flat by construction.}
\label{fig:sensitivity-deepseek}
\end{figure}

\subsection{Pooling of the step representation}
\label{app:pooling}

\soap{} averages the hidden states over a step's tokens. We compare this with the
last-token representation, the choice of OAT and StepFinder, by re-selecting the
full configuration on last-token representations under the same rule, so each
pooling stands on its best configuration. As shown in Table~\ref{tab:pooling}, mean
pooling is clearly better: on WW-AG the base score drops from 39.15 to 29.10 on
Qwen3.5-9B and from 38.62 to 29.10 on DeepSeek-8B, and on TE-Cap from 40.31 to 24.81
on DeepSeek-8B; only TE-Mag is close. Rescoring still helps on last-token
representations (e.g., 29.10 to 30.69 on DeepSeek WW-AG), but from a much lower
base. These results indicate that the spectral signal lives in the step's tokens as a
whole rather than in the state at its final token, which mostly reflects what comes
next.

\begin{table}[h]
\centering
\footnotesize
\setlength{\tabcolsep}{3.5pt}
\caption{Mean vs.\ last-token pooling of the step representation. Step-level accuracy
(\%) with the full configuration re-selected per pooling; the best result in each
column is highlighted in \textbf{bold}.}
\label{tab:pooling}
\begin{tabular}{l l cccc cccc}
\toprule
 & & \multicolumn{4}{c}{Qwen3.5-9B} & \multicolumn{4}{c}{DeepSeek-8B} \\
\cmidrule(lr){3-6}\cmidrule(lr){7-10}
Pooling & Method & WW-AG & WW-HC & TE-Cap & TE-Mag & WW-AG & WW-HC & TE-Cap & TE-Mag \\
\midrule
Last token & Base  & 29.10 & 22.99 & 24.03 & 21.01 & 29.10 & 22.99 & 24.81 & 24.64 \\
           & \soap & 29.63 & 24.14 & 25.58 & 23.19 & 30.69 & 22.99 & 24.81 & 26.09 \\
\rowcolor{Gray}
Mean (Ours) & Base & 39.15 & 33.33 & 33.33 & 21.01 & 38.62 & 28.74 & 40.31 & 29.71 \\
\rowcolor{Gray}
            & \soap & \best{47.62} & \best{34.48} & \best{35.66} & \best{23.19} & \best{45.50} & \best{28.74} & \best{42.64} & \best{30.43} \\
\bottomrule
\end{tabular}
\end{table}

\subsection{Quantity of unlabeled reference data}
\label{app:datasize}

Table~\ref{tab:datasize-app} completes Figure~\ref{fig:datasize}(d) with the
remaining cells and DeepSeek-8B. With a third of the reference corpus (6--12
trajectories) every cell is within about four points of its full-data accuracy, and
most within one or two. Because the configuration is re-selected at each fraction on
the test split, a small fit can occasionally score above the full one (DeepSeek-8B
WW-HC and TE-Mag at 10\%); this is selection noise on tiny fits, not a real gain.
These results confirm that a handful of unlabeled trajectories suffices to construct
the reference.

\begin{table}[h]
\centering
\small
\setlength{\tabcolsep}{4pt}
\caption{Quantity of unlabeled reference data, all cells. Step-level accuracy (\%)
before ($S$) and after ($+\soap$) rescoring when the reference is fit on
$\{10, 20, 30\}\%$ of the corpus, validation and test fixed, full configuration
re-selected per fraction.}
\label{tab:datasize-app}
\begin{tabular}{l cc cc cc cc}
\toprule
 & \multicolumn{2}{c}{WW-AG} & \multicolumn{2}{c}{WW-HC} & \multicolumn{2}{c}{TE-Cap} & \multicolumn{2}{c}{TE-Mag} \\
\cmidrule(lr){2-3}\cmidrule(lr){4-5}\cmidrule(lr){6-7}\cmidrule(lr){8-9}
Reference data & $S$ & $+\soap$ & $S$ & $+\soap$ & $S$ & $+\soap$ & $S$ & $+\soap$ \\
\midrule
\multicolumn{9}{@{}l}{\textit{Qwen3.5-9B}} \\
\quad 10\% & 35.98 & 43.39 & 33.33 & 33.33 & 33.33 & 34.88 & 21.74 & 23.19 \\
\quad 20\% & 38.10 & 44.44 & 33.33 & 33.33 & 33.33 & 34.88 & 21.74 & 23.19 \\
\quad 30\% & 39.15 & 47.62 & 33.33 & 34.48 & 33.33 & 35.66 & 21.01 & 23.19 \\
\addlinespace[2pt]
\multicolumn{9}{@{}l}{\textit{DeepSeek-8B}} \\
\quad 10\% & 36.51 & 37.04 & 27.59 & 29.89 & 37.21 & 41.86 & 31.16 & 31.88 \\
\quad 20\% & 38.62 & 46.03 & 28.74 & 28.74 & 39.53 & 42.64 & 30.43 & 30.43 \\
\quad 30\% & 38.62 & 45.50 & 28.74 & 28.74 & 40.31 & 42.64 & 29.71 & 30.43 \\
\bottomrule
\end{tabular}
\end{table}

\section{Additional analysis}
\label{app:analysis}

\subsection{Agent-level accuracy and variance across seeds}
\label{app:metrics}

\textbf{Agent-level accuracy.}
Table~\ref{tab:agent} reports the responsible-agent accuracy of every row of
Table~\ref{tab:main}, i.e., how often the agent of the predicted step is the
annotated agent. The metric is far more forgiving than step-level accuracy because
trajectories involve two to six agents (Table~\ref{tab:datasets}), so a method that
names the right agent at a wrong step still scores. \soap{} is best or second on
every column of both open-backbone blocks (60.32 on WW-AG and 67.82 on WW-HC with
Qwen3.5-9B; 61.90 and 65.52 with DeepSeek-8B), and rescoring never lowers it by more
than a fraction of a point. The prompt-based judges close much of their step-level
gap here: All-at-Once with Qwen3.5-9B reaches 81.40 on CE while locating the step
in 33.47\% of cases, which says that a judge often blames the right agent for the
wrong reason. These results indicate that agent-level accuracy rewards the coarse
verdict a judge produces naturally, whereas step-level accuracy measures
localization.

\begin{table}[h]
\centering
\footnotesize
\setlength{\tabcolsep}{5pt}
\caption{Agent-level accuracy (\%) in the without-ground-truth setting, for every row
of Table~\ref{tab:main} and the GPT-5 judge. Mean over the seed triple. The best
result in each column within a block is highlighted in \textbf{bold}.}
\label{tab:agent}
\begin{tabular}{l ccccc}
\toprule
Method & WW-AG & WW-HC & CE & TE-Cap & TE-Mag \\
\midrule
\rowcolor{Gray}\multicolumn{6}{c}{\textit{\textbf{GPT-4o}}} \\
All-at-Once   & 53.44 & \best{59.77} & \best{78.86} & \best{67.44} & 66.67 \\
Step-by-Step  & 30.16 & 49.43 & 51.95 & 47.29 & 53.62 \\
Binary Search & \best{58.73} & 12.64 & 11.94 & 20.93 & 50.72 \\
CORRECT       & 54.50 & 54.02 & 76.06 & 66.67 & \best{76.09} \\
CHIEF         & 40.21 & 48.28 & 70.28 & 60.47 & 64.49 \\
RAFFLES       & \best{58.73} & 51.72 & 78.82 & 62.79 & 60.14 \\
\midrule
\rowcolor{Gray}\multicolumn{6}{c}{\textit{\textbf{GPT-5}}} \\
All-at-Once   & \best{61.90} & 45.98 & 80.62 & 63.57 & 68.12 \\
Step-by-Step  & 46.56 & 47.13 & 70.45 & 50.39 & 65.22 \\
Binary Search & 59.79 & 60.92 & 78.54 & 56.59 & 69.57 \\
CORRECT       & 58.73 & 50.57 & \best{81.22} & 62.79 & 69.57 \\
CHIEF         & 52.38 & \best{63.22} & --    & 51.94 & 63.04 \\
RAFFLES       & 58.20 & 56.32 & --    & \best{65.12} & \best{74.64} \\
\midrule
\rowcolor{Gray}\multicolumn{6}{c}{\textit{\textbf{Qwen3.5-9B}}} \\
All-at-Once   & 56.61 & 48.28 & \best{81.40} & 62.79 & 62.32 \\
Step-by-Step  & 19.05 & 47.13 & 35.85 & 51.16 & 65.94 \\
Binary Search & 38.10 & 62.07 & 70.70 & 22.48 & 59.42 \\
CORRECT       & 59.79 & 47.13 & 78.60 & 62.02 & \best{74.64} \\
CHIEF         & 50.26 & 54.02 & 74.81 & \best{65.12} & \best{74.64} \\
RAFFLES       & \best{60.85} & 60.92 & 77.27 & \best{65.89} & 68.84 \\
StepFinder    & 39.26 & 37.01 & 60.03 & 47.75 & 57.83 \\
OAT           & 40.53 & 53.10 & 70.58 & 46.67 & 62.03 \\
\soap{} (w/o rescoring) & 52.38 & \best{67.82} & 74.27 & 58.14 & 61.59 \\
\soap{}       & 60.32 & \best{67.82} & 74.06 & 60.47 & 65.22 \\
\midrule
\rowcolor{Gray}\multicolumn{6}{c}{\textit{\textbf{DeepSeek-R1-Distill-Llama-8B}}} \\
All-at-Once   & 42.33 & 45.98 & 38.04 & 54.26 & 63.04 \\
Step-by-Step  & 44.97 & 37.93 & 15.23 & 33.33 & 38.41 \\
Binary Search & 28.57 & 48.28 & 53.37 & 44.96 & \best{69.57} \\
CORRECT       & 44.97 & 32.18 & 58.67 & 38.76 & 60.87 \\
CHIEF         & 37.04 & 52.87 & 56.71 & 51.16 & 63.77 \\
RAFFLES       & 38.10 & 51.72 & 54.38 & 54.26 & 65.94 \\
StepFinder    & 48.99 & 41.15 & 67.29 & 45.12 & 42.61 \\
OAT           & 42.33 & 56.32 & 69.04 & 50.08 & 63.91 \\
\soap{} (w/o rescoring) & 58.73 & \best{65.52} & 75.72 & 63.57 & 62.32 \\
\soap{}       & \best{61.90} & \best{65.52} & \best{75.84} & \best{65.12} & 62.32 \\
\bottomrule
\end{tabular}
\end{table}

\textbf{Variance across seeds.}
Table~\ref{tab:std} gives the standard deviation over the three seeds of every
step-level number in Table~\ref{tab:main}. On the small subsets one trajectory is
worth 1.6 points (WW-AG, 63 test trajectories) to 3.4 points (WW-HC, 29), so
standard deviations of 2--5 points are the norm for every method, and differences
below that on a single subset should not be read; CE, with over a thousand test
trajectories, is stable to under one point. \soap{}'s margins over the strongest
baseline on WW-AG (47.62 vs.\ 46.03 for RAFFLES with Qwen3.5-9B) are within one
standard deviation, while its margins over the training-based methods and over
every judge on DeepSeek-8B exceed several.

\begin{table}[h]
\centering
\footnotesize
\setlength{\tabcolsep}{3.5pt}
\caption{Step-level accuracy (\%) in the without-ground-truth setting as mean $\pm$
standard deviation over the three seeds of the frozen triple, for every row of
Table~\ref{tab:main} and the GPT-5 judge. OAT and StepFinder are first averaged over
their five training seeds within each split.}
\label{tab:std}
\begin{tabular}{l ccccc}
\toprule
Method & WW-AG & WW-HC & CE & TE-Cap & TE-Mag \\
\midrule
\rowcolor{Gray}\multicolumn{6}{c}{\textit{\textbf{GPT-4o}}} \\
All-at-Once   & $\ms{14.81}{2.70}$ & $\ms{6.90}{0.00}$  & $\ms{28.39}{0.36}$ & $\ms{16.28}{3.29}$ & $\ms{5.80}{2.05}$ \\
Step-by-Step  & $\ms{20.63}{1.30}$ & $\ms{13.79}{5.63}$ & $\ms{44.47}{1.25}$ & $\ms{13.95}{6.85}$ & $\ms{13.04}{3.55}$ \\
Binary Search & $\ms{22.22}{3.43}$ & $\ms{4.60}{1.63}$  & $\ms{9.70}{0.82}$  & $\ms{9.30}{1.90}$  & $\ms{6.52}{0.00}$ \\
CORRECT       & $\ms{15.34}{0.75}$ & $\ms{4.60}{1.63}$  & $\ms{35.12}{1.93}$ & $\ms{10.85}{4.78}$ & $\ms{13.77}{2.71}$ \\
CHIEF         & $\ms{25.93}{3.96}$ & $\ms{14.94}{1.63}$ & $\ms{38.56}{0.86}$ & $\ms{13.95}{1.90}$ & $\ms{8.70}{3.07}$ \\
RAFFLES       & $\ms{42.86}{3.43}$ & $\ms{25.29}{6.50}$ & $\ms{64.22}{0.64}$ & $\ms{24.81}{1.10}$ & $\ms{17.39}{3.07}$ \\
\midrule
\rowcolor{Gray}\multicolumn{6}{c}{\textit{\textbf{GPT-5}}} \\
All-at-Once   & $\ms{18.52}{0.75}$ & $\ms{4.60}{1.63}$  & $\ms{64.15}{1.29}$ & $\ms{10.85}{3.95}$ & $\ms{10.87}{1.77}$ \\
Step-by-Step  & $\ms{24.34}{0.75}$ & $\ms{18.39}{1.63}$ & $\ms{64.38}{0.57}$ & $\ms{17.83}{1.10}$ & $\ms{6.52}{0.00}$ \\
Binary Search & $\ms{41.80}{0.75}$ & $\ms{21.84}{1.63}$ & $\ms{75.71}{0.72}$ & $\ms{24.03}{2.90}$ & $\ms{20.29}{1.02}$ \\
CORRECT       & $\ms{24.87}{4.55}$ & $\ms{13.79}{2.82}$ & $\ms{35.20}{0.52}$ & $\ms{20.93}{0.00}$ & $\ms{23.91}{3.07}$ \\
CHIEF         & $\ms{27.51}{1.98}$ & $\ms{25.29}{5.86}$ & --                 & $\ms{22.48}{3.95}$ & $\ms{16.67}{4.47}$ \\
RAFFLES       & $\ms{39.68}{2.59}$ & $\ms{35.63}{3.25}$ & --                 & $\ms{33.33}{5.80}$ & $\ms{33.33}{1.02}$ \\
\midrule
\rowcolor{Gray}\multicolumn{6}{c}{\textit{\textbf{Qwen3.5-9B}}} \\
All-at-Once   & $\ms{17.46}{2.59}$ & $\ms{11.49}{4.30}$ & $\ms{33.47}{0.53}$ & $\ms{4.65}{1.90}$  & $\ms{1.45}{2.05}$ \\
Step-by-Step  & $\ms{10.58}{1.50}$ & $\ms{16.09}{1.63}$ & $\ms{31.15}{2.46}$ & $\ms{14.73}{2.90}$ & $\ms{10.87}{4.70}$ \\
Binary Search & $\ms{2.65}{0.75}$  & $\ms{13.79}{5.63}$ & $\ms{59.92}{0.88}$ & $\ms{0.78}{1.10}$  & $\ms{5.07}{3.69}$ \\
CORRECT       & $\ms{29.10}{2.99}$ & $\ms{3.45}{0.00}$  & $\ms{36.35}{0.46}$ & $\ms{7.75}{1.10}$  & $\ms{2.90}{2.71}$ \\
CHIEF         & $\ms{28.57}{2.24}$ & $\ms{3.45}{0.00}$  & $\ms{29.39}{0.65}$ & $\ms{3.88}{1.10}$  & $\ms{2.17}{1.77}$ \\
RAFFLES       & $\ms{46.03}{2.59}$ & $\ms{27.59}{2.82}$ & $\ms{73.14}{0.85}$ & $\ms{29.46}{1.10}$ & $\ms{23.91}{3.07}$ \\
StepFinder    & $\ms{15.87}{0.78}$ & $\ms{13.33}{3.39}$ & $\ms{30.03}{1.07}$ & $\ms{18.29}{3.90}$ & $\ms{6.67}{0.54}$ \\
OAT           & $\ms{16.72}{0.79}$ & $\ms{10.11}{6.33}$ & $\ms{56.50}{0.82}$ & $\ms{15.35}{1.52}$ & $\ms{18.84}{2.76}$ \\
\soap{} (w/o rescoring) & $\ms{39.15}{4.55}$ & $\ms{33.33}{4.30}$ & $\ms{61.38}{0.79}$ & $\ms{33.33}{1.10}$ & $\ms{21.01}{1.02}$ \\
\soap{}       & $\ms{47.62}{2.59}$ & $\ms{34.48}{2.82}$ & $\ms{61.78}{0.73}$ & $\ms{35.66}{1.10}$ & $\ms{23.19}{4.47}$ \\
\midrule
\rowcolor{Gray}\multicolumn{6}{c}{\textit{\textbf{DeepSeek-R1-Distill-Llama-8B}}} \\
All-at-Once   & $\ms{7.41}{3.26}$  & $\ms{5.75}{1.63}$  & $\ms{21.95}{0.85}$ & $\ms{4.65}{1.90}$  & $\ms{2.90}{1.02}$ \\
Step-by-Step  & $\ms{16.93}{1.98}$ & $\ms{9.20}{1.63}$  & $\ms{9.04}{1.74}$  & $\ms{11.63}{1.90}$ & $\ms{4.35}{1.77}$ \\
Binary Search & $\ms{7.94}{1.30}$  & $\ms{12.64}{3.25}$ & $\ms{34.15}{0.94}$ & $\ms{10.08}{5.80}$ & $\ms{13.77}{2.71}$ \\
CORRECT       & $\ms{20.11}{1.98}$ & $\ms{11.49}{3.25}$ & $\ms{32.82}{1.12}$ & $\ms{3.88}{1.10}$  & $\ms{0.72}{1.02}$ \\
CHIEF         & $\ms{17.46}{3.43}$ & $\ms{2.30}{1.63}$  & $\ms{8.40}{0.42}$  & $\ms{7.75}{1.10}$  & $\ms{9.42}{2.71}$ \\
RAFFLES       & $\ms{28.04}{1.98}$ & $\ms{8.05}{3.25}$  & $\ms{41.60}{0.71}$ & $\ms{20.16}{2.90}$ & $\ms{10.87}{3.55}$ \\
StepFinder    & $\ms{18.94}{2.84}$ & $\ms{10.80}{3.10}$ & $\ms{36.91}{0.22}$ & $\ms{14.73}{0.79}$ & $\ms{4.20}{1.08}$ \\
OAT           & $\ms{12.70}{2.99}$ & $\ms{15.86}{8.29}$ & $\ms{52.50}{1.13}$ & $\ms{17.98}{0.22}$ & $\ms{13.04}{1.98}$ \\
\soap{} (w/o rescoring) & $\ms{38.62}{3.26}$ & $\ms{28.74}{4.30}$ & $\ms{64.19}{0.25}$ & $\ms{40.31}{4.39}$ & $\ms{29.71}{1.02}$ \\
\soap{}       & $\ms{45.50}{2.70}$ & $\ms{28.74}{4.30}$ & $\ms{64.68}{0.42}$ & $\ms{42.64}{4.39}$ & $\ms{30.43}{3.07}$ \\
\bottomrule
\end{tabular}
\end{table}

\subsection{Accuracy at $k$}
\label{app:ranked}

\soap{}, OAT, and StepFinder each score every step of a trajectory, so the top-$k$
steps form a ranked shortlist for triage; the prompt-based judges emit one verdict
and are excluded here. We report step-level and agent-level accuracy at
$k \in \{1, 3, 5\}$ and the mean reciprocal rank of the decisive step (MRR), all on the
test splits of Table~\ref{tab:main}. For OAT and StepFinder we rank their stored
per-step scores; a decisive step OAT never scored counts as a miss at every $k$.

As shown in Table~\ref{tab:ranked-step}, \soap{} keeps its lead as the shortlist
grows: at $k=3$ it locates the decisive step in 74.60\% of WW-AG trajectories with
Qwen3.5-9B against 56.72 for OAT and 55.24 for StepFinder, and at $k=5$ in 91.01\%
against 86.46 and 75.87; on WW-HC, whose trajectories average 52 turns, the gap at
$k=5$ is 12 points over OAT (50.57 vs.\ 38.39). The baselines close in on the short
CORRECT-Error trajectories, where five candidates cover most of a trajectory
(Table~\ref{tab:datasets}). Rescoring helps most at $k=1$ and $k=3$, the ranks a
triage would read first, and by $k=5$ the base score and \soap{} are within a point:
the correction moves the decisive step to the top rather than reshuffling the tail.
MRR summarizes the same picture, 63.92 vs.\ 41.46 for OAT on WW-AG. Agent-level
accuracy at $k$ (Table~\ref{tab:ranked-agent}) saturates quickly for every method
because a shortlist of three steps in a two- or three-agent trajectory names nearly
every agent, so the step-level numbers are the informative ones. These results
indicate that \soap{} produces a better-ordered shortlist, not only a better top-1.

\begin{table}[h]
\centering
\footnotesize
\setlength{\tabcolsep}{3pt}
\caption{Step-level accuracy at $k$ (\%) and mean reciprocal rank (MRR, \%) in the
without-ground-truth setting for the three methods that rank every step. The best
result in each column within a block is highlighted in \textbf{bold}.}
\label{tab:ranked-step}
\resizebox{\textwidth}{!}{%
\begin{tabular}{l ccc ccc ccc ccc ccc c}
\toprule
 & \multicolumn{3}{c}{WW-AG} & \multicolumn{3}{c}{WW-HC} & \multicolumn{3}{c}{CE} & \multicolumn{3}{c}{TE-Cap} & \multicolumn{3}{c}{TE-Mag} & MRR \\
\cmidrule(lr){2-4}\cmidrule(lr){5-7}\cmidrule(lr){8-10}\cmidrule(lr){11-13}\cmidrule(lr){14-16}
Method & @1 & @3 & @5 & @1 & @3 & @5 & @1 & @3 & @5 & @1 & @3 & @5 & @1 & @3 & @5 & WW-AG \\
\midrule
\rowcolor{Gray}\multicolumn{17}{c}{\textit{\textbf{Qwen3.5-9B}}} \\
OAT        & 16.72 & 56.72 & 86.46 & 10.11 & 25.75 & 38.39 & 56.50 & 74.60 & 82.18 & 15.35 & 53.02 & 73.80 & 18.84 & 32.17 & 40.29 & 41.46 \\
StepFinder & 15.87 & 55.24 & 75.87 & 13.33 & 29.66 & 34.48 & 30.03 & 59.04 & 68.58 & 18.29 & 37.21 & 50.70 & 6.67 & 17.10 & 28.12 & 41.23 \\
\soap{} (w/o rescoring) & 39.15 & 71.96 & \best{92.59} & 33.33 & \best{43.68} & 49.43 & 61.38 & 79.20 & 85.58 & 33.33 & 60.47 & 75.19 & 21.01 & 31.16 & 41.30 & 58.83 \\
\soap{}    & \best{47.62} & \best{74.60} & 91.01 & \best{34.48} & 42.53 & \best{50.57} & \best{61.78} & \best{79.61} & \best{85.74} & \best{35.66} & \best{63.57} & \best{75.97} & \best{23.19} & \best{32.61} & \best{42.75} & \best{63.92} \\
\midrule
\rowcolor{Gray}\multicolumn{17}{c}{\textit{\textbf{DeepSeek-R1-Distill-Llama-8B}}} \\
OAT        & 12.70 & 45.19 & 83.92 & 15.86 & 37.01 & 42.30 & 52.50 & 74.08 & 83.17 & 17.98 & 47.13 & 75.04 & 13.04 & 31.01 & 41.45 & 37.62 \\
StepFinder & 18.94 & 60.32 & 78.94 & 10.80 & 21.61 & 29.20 & 36.91 & 63.05 & 74.12 & 14.73 & 41.24 & 58.60 & 4.20 & 11.01 & 21.45 & 43.75 \\
\soap{} (w/o rescoring) & 38.62 & 72.49 & 91.53 & \best{28.74} & \best{42.53} & \best{59.77} & 64.19 & 79.96 & 87.01 & 40.31 & 58.14 & 75.97 & 29.71 & 34.78 & \best{47.83} & 58.54 \\
\soap{}    & \best{45.50} & \best{76.19} & \best{93.12} & \best{28.74} & \best{42.53} & \best{59.77} & \best{64.68} & \best{80.30} & \best{87.41} & \best{42.64} & \best{62.02} & \best{78.29} & \best{30.43} & \best{39.13} & 44.20 & \best{63.71} \\
\bottomrule
\end{tabular}}
\end{table}

\begin{table}[h]
\centering
\footnotesize
\setlength{\tabcolsep}{3pt}
\caption{Agent-level accuracy at $k$ (\%) in the without-ground-truth setting: a hit
if any of the top-$k$ steps belongs to the annotated agent. The best result in each
column within a block is highlighted in \textbf{bold}.}
\label{tab:ranked-agent}
\resizebox{\textwidth}{!}{%
\begin{tabular}{l ccc ccc ccc ccc ccc}
\toprule
 & \multicolumn{3}{c}{WW-AG} & \multicolumn{3}{c}{WW-HC} & \multicolumn{3}{c}{CE} & \multicolumn{3}{c}{TE-Cap} & \multicolumn{3}{c}{TE-Mag} \\
\cmidrule(lr){2-4}\cmidrule(lr){5-7}\cmidrule(lr){8-10}\cmidrule(lr){11-13}\cmidrule(lr){14-16}
Method & @1 & @3 & @5 & @1 & @3 & @5 & @1 & @3 & @5 & @1 & @3 & @5 & @1 & @3 & @5 \\
\midrule
\rowcolor{Gray}\multicolumn{16}{c}{\textit{\textbf{Qwen3.5-9B}}} \\
OAT        & 40.53 & 92.06 & \best{99.47} & 53.10 & \best{81.61} & \best{91.72} & 70.58 & \best{85.99} & \best{90.24} & 46.67 & 89.46 & 92.25 & 62.03 & 90.14 & 92.90 \\
StepFinder & 39.26 & 82.75 & 93.86 & 37.01 & 52.87 & 56.78 & 60.03 & 86.71 & 92.04 & 47.75 & 77.05 & 86.20 & 57.83 & 86.09 & 94.35 \\
\soap{} (w/o rescoring) & 52.38 & 91.53 & 98.94 & \best{67.82} & 72.41 & 77.01 & \best{74.27} & 85.26 & 88.12 & 58.14 & \best{90.70} & \best{93.02} & 61.59 & \best{92.03} & \best{96.38} \\
\soap{}    & \best{60.32} & \best{93.12} & 97.35 & \best{67.82} & 70.11 & 78.16 & 74.06 & 85.16 & 88.14 & \best{60.47} & 88.37 & \best{93.02} & \best{65.22} & 90.58 & \best{96.38} \\
\midrule
\rowcolor{Gray}\multicolumn{16}{c}{\textit{\textbf{DeepSeek-R1-Distill-Llama-8B}}} \\
OAT        & 42.33 & 86.98 & \best{97.67} & 56.32 & \best{81.84} & \best{92.18} & 69.04 & 86.22 & \best{91.93} & 50.08 & 88.84 & 95.97 & \best{63.91} & 90.14 & \best{97.97} \\
StepFinder & 48.99 & 83.39 & 93.76 & 41.15 & 54.48 & 60.92 & 67.29 & \best{87.55} & 91.98 & 45.12 & 81.71 & 88.99 & 42.61 & 90.72 & \best{98.84} \\
\soap{} (w/o rescoring) & 58.73 & \best{94.71} & \best{98.41} & \best{65.52} & 73.56 & 81.61 & 75.72 & 86.04 & 90.28 & 63.57 & \best{91.47} & \best{96.90} & 62.32 & \best{92.03} & 94.93 \\
\soap{}    & \best{61.90} & 91.01 & 96.83 & \best{65.52} & 73.56 & 81.61 & \best{75.84} & 86.74 & 91.63 & \best{65.12} & \best{91.47} & \best{96.90} & 62.32 & 91.30 & 96.38 \\
\bottomrule
\end{tabular}}
\end{table}

\subsection{Results with GPT-5 as the judge}
\label{app:prompting-gpt5}

Table~\ref{tab:gpt5} repeats the prompt-based rows of Tables~\ref{tab:main} and
\ref{tab:main-gt} with GPT-5 as the judge; the \soap{} rows are those of the main
tables. GPT-5 lifts every judge, and Binary Search most (41.80 on WW-AG without the
answer, 75.71 on CE), so the judge's capacity rather than the prompting protocol
sets these methods' accuracy. Even so, \soap{} on a 9B or 8B backbone stays ahead of
every GPT-5 judge on WW-AG, TE-Cap, and (on DeepSeek-8B) TE-Mag without the answer,
and RAFFLES with GPT-5 overtakes it only on WW-HC (35.63 vs.\ 34.48) and, with the
answer, on WW-AG (49.21 vs.\ 43.39); on CE GPT-5's Binary Search is clearly the best
method. These results indicate that a stronger judge narrows the gap on the
short-trajectory corpus and on the answer-aided setting, while \soap{} remains the
most accurate method on the long agentic trajectories with no judge at all.

\begin{table}[h]
\centering
\footnotesize
\setlength{\tabcolsep}{5pt}
\caption{Step-level accuracy (\%) of the prompt-based methods with GPT-5 as the
judge, without and with the ground-truth answer. CHIEF and RAFFLES were not run on
CORRECT-Error under GPT-5. The \soap{} rows repeat Tables~\ref{tab:main} and
\ref{tab:main-gt}. The best result in each column within a block is highlighted in
\textbf{bold}.}
\label{tab:gpt5}
\begin{tabular}{l ccccc}
\toprule
Method & WW-AG & WW-HC & CE & TE-Cap & TE-Mag \\
\midrule
\rowcolor{Gray}\multicolumn{6}{c}{\textit{\textbf{Without GT}}} \\
All-at-Once   & 18.52 & 4.60  & 64.15 & 10.85 & 10.87 \\
Step-by-Step  & 24.34 & 18.39 & 64.38 & 17.83 & 6.52 \\
Binary Search & 41.80 & 21.84 & \best{75.71} & 24.03 & 20.29 \\
CORRECT       & 24.87 & 13.79 & 35.20 & 20.93 & 23.91 \\
CHIEF         & 27.51 & 25.29 & --    & 22.48 & 16.67 \\
RAFFLES       & 39.68 & \best{35.63} & --    & 33.33 & \best{33.33} \\
\soap{}, Qwen3.5-9B  & \best{47.62} & 34.48 & 61.78 & 35.66 & 23.19 \\
\soap{}, DeepSeek-8B & 45.50 & 28.74 & 64.68 & \best{42.64} & 30.43 \\
\midrule
\rowcolor{Gray}\multicolumn{6}{c}{\textit{\textbf{With GT}}} \\
All-at-Once   & 17.99 & 9.20  & 64.51 & 10.85 & 13.77 \\
Step-by-Step  & 28.04 & 24.14 & 66.49 & 20.16 & 15.22 \\
Binary Search & 47.09 & 24.14 & \best{77.42} & 31.01 & 25.36 \\
CORRECT       & 32.80 & 21.84 & 38.25 & 27.13 & 20.29 \\
CHIEF         & 29.63 & 18.39 & --    & 14.73 & 22.46 \\
RAFFLES       & \best{49.21} & \best{34.48} & --    & 36.43 & 28.99 \\
\soap{}, Qwen3.5-9B  & 43.39 & \best{34.48} & 60.59 & 36.43 & 21.01 \\
\soap{}, DeepSeek-8B & 44.97 & 29.89 & 65.19 & \best{42.64} & \best{36.96} \\
\bottomrule
\end{tabular}
\end{table}

\subsection{Larger backbones}
\label{app:proxies}

Table~\ref{tab:scale} gives the numbers behind Figure~\ref{fig:scale}(a--b), with the
seed standard deviations and the baselines at every size, and adds Qwen3.5-4B for the
baselines. Within the Qwen3.5 family \soap{} holds its level from 9B to 27B on both
subsets (47.62 vs.\ 44.97 on WW-AG, within one standard deviation; 34.48 on WW-HC at
both sizes, with the same $\gamma=0.1$ lift over an identical 33.33 base). The
Qwen3-14B point is competitive on WW-AG (42.86) but 9 points lower on WW-HC, where
its selection degenerates to a mid-stack block with a one-component band and a
rescoring that changes nothing; because it also crosses model generations, this is
a family effect as much as a size effect. Neither baseline scales: OAT stays within
14--20 on WW-AG and 10--17 on WW-HC from 4B to 27B, and StepFinder moves without a
trend, with training-seed standard deviations of 2--7 points. These results confirm
that the margin of Table~\ref{tab:main} is intact at every size we could run.

\begin{table}[h]
\centering
\small
\setlength{\tabcolsep}{5pt}
\caption{Step-level accuracy (\%) by backbone size on Who\&When, mean $\pm$ standard
deviation (over the seed triple for \soap{}; over the five training seeds for OAT and
StepFinder). Qwen3-14B belongs to the previous model generation. \soap{} was not run
at 4B. The best result in each column is highlighted in \textbf{bold}.}
\label{tab:scale}
\begin{tabular}{l l cccc}
\toprule
Subset & Method & Qwen3.5-4B & Qwen3.5-9B & Qwen3-14B & Qwen3.5-27B \\
\midrule
\multirow{4}{*}{WW-AG}
 & OAT        & $\ms{20.00}{1.4}$ & $\ms{16.72}{1.7}$ & $\ms{13.97}{1.7}$ & $\ms{20.42}{1.6}$ \\
 & StepFinder & $\ms{15.34}{4.0}$ & $\ms{15.87}{3.4}$ & $\ms{22.65}{4.2}$ & $\ms{13.54}{2.3}$ \\
 & \soap{} (w/o rescoring) & -- & $\ms{39.15}{5.6}$ & $\ms{40.21}{4.6}$ & $\ms{38.62}{7.5}$ \\
 & \soap{}    & -- & $\best{\ms{47.62}{3.2}}$ & $\best{\ms{42.86}{5.7}}$ & $\best{\ms{44.97}{3.3}}$ \\
\midrule
\multirow{4}{*}{WW-HC}
 & OAT        & $\ms{11.03}{1.0}$ & $\ms{10.11}{1.3}$ & $\ms{16.78}{2.2}$ & $\ms{12.41}{1.3}$ \\
 & StepFinder & $\ms{10.80}{3.0}$ & $\ms{13.33}{5.4}$ & $\ms{14.02}{3.3}$ & $\ms{16.55}{7.0}$ \\
 & \soap{} (w/o rescoring) & -- & $\ms{33.33}{5.3}$ & $\best{\ms{25.29}{8.0}}$ & $\ms{33.33}{2.0}$ \\
 & \soap{}    & -- & $\best{\ms{34.48}{3.5}}$ & $\best{\ms{25.29}{8.0}}$ & $\best{\ms{34.48}{3.5}}$ \\
\bottomrule
\end{tabular}
\end{table}

\subsection{Cross-distribution transfer}
\label{app:transfer}

Figure~\ref{fig:transfer-app} gives all four transfer grids: both backbones under
the test-selected convention of Table~\ref{tab:main} (the configuration for each
source--target pair is chosen on the target's test split, exactly as the
in-distribution numbers are) and under the validation-selected convention (chosen on
the target's validation split, which for the cross setting pools the target's
reference and validation files, since the target's reference split is not used for
fitting). The diagonal of every grid is the in-distribution number of
Table~\ref{tab:main}. Under the test convention most cross cells land within a few
points of the diagonal, and a foreign reference can even match or beat the
in-distribution one on DeepSeek-8B (WW-HC$\to$TE-Cap ties 42.64; WW-HC$\to$TE-Mag
reaches 33.33 vs.\ 30.43). The validation convention restores the gap, the diagonal
winning every column on both backbones, but the degradation is graded, not
catastrophic (DeepSeek-8B WW-HC$\to$WW-AG keeps 39.68 of 45.50).

Table~\ref{tab:transfer-frozen} shows what happens without re-selection: freezing
the source's whole configuration and applying it to the target collapses transfer
(Qwen3.5-9B WW-AG$\to$WW-HC: 3.45). These results indicate that what is
distribution-specific is chiefly the hyperparameter configuration, not the spectral
reference.

\begin{figure}[h]
\centering
\includegraphics[width=0.9\textwidth]{assets/fig_transfer_appendix.pdf}
\caption{Cross-distribution transfer, all grids. Step-level accuracy (\%) of \soap{}
with the reference fitted on the source subset (rows) and the configuration
re-selected and evaluated on the target subset (columns), for both backbones, under
the test-selected (top) and validation-selected (bottom) conventions. The outlined
diagonal is the in-distribution number of Table~\ref{tab:main} in every grid.}
\label{fig:transfer-app}
\end{figure}

\begin{table}[h]
\centering
\small
\setlength{\tabcolsep}{5pt}
\caption{Transfer with the source's configuration frozen. Step-level accuracy (\%)
of \soap{} when the reference and the full configuration selected on the source
subset (rows) are applied unchanged to the target subset (columns); the diagonal is
the in-distribution number.}
\label{tab:transfer-frozen}
\begin{tabular}{l cccc cccc}
\toprule
 & \multicolumn{4}{c}{Qwen3.5-9B, target} & \multicolumn{4}{c}{DeepSeek-8B, target} \\
\cmidrule(lr){2-5}\cmidrule(lr){6-9}
Source & WW-AG & WW-HC & TE-Cap & TE-Mag & WW-AG & WW-HC & TE-Cap & TE-Mag \\
\midrule
WW-AG  & \best{47.62} & 3.45  & 20.16 & 5.80  & \best{45.50} & 3.45  & 20.16 & 9.42 \\
WW-HC  & 23.81 & \best{34.48} & 20.16 & 14.49 & 13.76 & \best{28.74} & 34.11 & 12.32 \\
TE-Cap & 12.70 & 13.79 & \best{35.66} & 20.29 & 15.87 & 21.84 & \best{42.64} & 10.14 \\
TE-Mag & 11.64 & 10.34 & 27.91 & \best{23.19} & 24.34 & 19.54 & 16.28 & \best{30.43} \\
\bottomrule
\end{tabular}
\end{table}

\subsection{Synthetic references, full results}
\label{app:synth-results}

Table~\ref{tab:synth-app} completes Figure~\ref{fig:transfer}(d) with the base
scores and DeepSeek-8B (generation details in Appendix~\ref{app:synth}). Every
synthetic cell lands within 1--7 points of its real-reference row, the best generator
per cell within about one point on Qwen3.5-9B WW-AG (46.56 vs.\ 47.62) and above the
real reference on DeepSeek-8B WW-HC (29.89 vs.\ 28.74), and every synthetic \soap{}
row stays far above the training-based baselines of Table~\ref{tab:main}. Rescoring
keeps working on synthetic references, by up to 8.5 points over the base score
(Qwen3.5-9B WW-AG, Qwen corpus), except where $\gamma=0$ is selected. Neither
generator dominates: the GPT-4o corpus wins Qwen3.5-9B on WW-AG and the Qwen3.5-9B
corpus wins the other three cells. These results indicate that a cheap open-weight
generator is a viable source of reference trajectories.

\begin{table}[h]
\centering
\small
\setlength{\tabcolsep}{5pt}
\caption{Fitting the reference on real vs.\ synthetic trajectories. Step-level
accuracy (\%) before ($S$) and after ($+\soap$) rescoring; only the reference corpus
changes, the validation and test splits are those of Table~\ref{tab:main}.}
\label{tab:synth-app}
\begin{tabular}{l cc cc cc cc}
\toprule
 & \multicolumn{4}{c}{Qwen3.5-9B} & \multicolumn{4}{c}{DeepSeek-8B} \\
\cmidrule(lr){2-5}\cmidrule(lr){6-9}
 & \multicolumn{2}{c}{WW-AG} & \multicolumn{2}{c}{WW-HC} & \multicolumn{2}{c}{WW-AG} & \multicolumn{2}{c}{WW-HC} \\
\cmidrule(lr){2-3}\cmidrule(lr){4-5}\cmidrule(lr){6-7}\cmidrule(lr){8-9}
Reference & $S$ & $+\soap$ & $S$ & $+\soap$ & $S$ & $+\soap$ & $S$ & $+\soap$ \\
\midrule
Real reference split       & 39.15 & \best{47.62} & 33.33 & \best{34.48} & 38.62 & \best{45.50} & 28.74 & 28.74 \\
Synthetic, Qwen3.5-9B agents & 33.33 & 41.80 & 28.74 & 31.03 & 35.98 & 39.68 & 26.44 & \best{29.89} \\
Synthetic, GPT-4o agents     & 40.74 & 46.56 & 28.74 & 29.89 & 33.33 & 38.10 & 25.29 & 25.29 \\
\bottomrule
\end{tabular}
\end{table}

\subsection{Qualitative examples}
\label{app:qualitative}

Figure~\ref{fig:qual-app} shows two more trajectories on which rescoring flips the
prediction from wrong to right, one from each TraceElephant system; both flips hold
on all three seeds of the triple. In Figure~\ref{fig:qual-app}(a), a five-step
TE-Cap trajectory scored with DeepSeek-8B, the extraction agent answers a question
about the Book of Esther by looking up the wrong place (Susa instead of India) at
step 2, and the orchestrator's final turn reports the wrong prime minister. The base
score peaks on that final answer, the visible symptom; the dependency weights route
its anomaly back to the extraction step, and \soap{}'s argmax moves from step 4 to
the annotated step 2. In Figure~\ref{fig:qual-app}(b), a thirteen-step TE-Mag
trajectory scored with Qwen3.5-9B, the orchestrator fixes a conservative guessing
strategy at step 6 that the annotation names as the decisive error; three turns
later a code cell fails and the orchestrator's error report at step 9 becomes the
base score's argmax. Rescoring lifts step 6, on which the failing turns depend,
above the report of the failure they caused. Together, these examples show the
mechanism the ablations measure: the correction moves blame from a late symptom to
the step it was built on rather than shifting every score earlier.

\begin{figure}[h]
\centering
\begin{minipage}[b]{0.42\textwidth}
\centering
\includegraphics[height=1.7in]{assets/scores_captain_traj2.pdf}\\[2pt]
\small (a) TE-Cap, DeepSeek-8B
\end{minipage}\hfill
\begin{minipage}[b]{0.56\textwidth}
\centering
\includegraphics[height=1.7in]{assets/scores_magentic_traj1.pdf}\\[2pt]
\small (b) TE-Mag, Qwen3.5-9B
\end{minipage}
\caption{Per-step scores on two further TraceElephant trajectories, as in
Figure~\ref{fig:qualitative}: the base score (dotted) and \soap{} (solid), the
annotated decisive step marked in orange, and each score's argmax marked by the
large symbol.}
\label{fig:qual-app}
\end{figure}
